% Appendix B.3.1: regularization of the closed-form TimesFM-3 adapter (Phase 2g, review period).
% Numbers: results/three_model_phase2g_timesfm3_zero_ridge/ (Narval job 9588346) and
% experiments/make_rebuttal_numbers.py, section D. Figure: experiments/make_rebuttal_figures.py.

\subsubsection{Regularization of the TimesFM-3 adapter}
\label{sec:app_timesfm3_ridge}

The submitted closed-form TimesFM-3 adapter penalized $\lambda\|\Delta_\ell\|_F^2$, i.e.\ it shrank $A_\ell=I+\Delta_\ell$ toward the identity and therefore toward the hard cut. For TimesFM-3 the hard cut is 35 to 160 times worse than the full model in median test MASE from L0 to L17 (Table~\ref{tab:timesfm3_ridge}). Validation therefore selected almost no regularization, $\kappa\le10^{-8}$ in 231 of the 280 fits (14 datasets $\times$ 20 depths), and the adapter nearly interpolated its 1,394 training windows: from L1 to L17, its median validation distillation error was 13 to 184 times its training error, and its test $R^2$ against the full model's forecast was negative on 4 to 11 of the 14 datasets (Table~\ref{tab:timesfm3_ridge_fit}). At L0 the same fit retained far fewer effective parameters and generalized, which is why aligned TimesFM-3 at L0 outperformed L2 to L17 in the submission.

\paragraph{Zero anchor.} We refit the adapter with the penalty on the whole linear map, $\lambda\|I+\Delta_\ell\|_F^2$ (Equation~\ref{eq:adapter_ridge}), which shrinks $A_\ell$ toward zero, and changed nothing else: the same cached states, splits, teacher forecasts, metric $W^\top W$, $\lambda=\kappa\,g_{\max}m_{\max}$ with $\kappa\in\{10^{-12},\dots,10^{2}\}$, and selection on validation distillation error, with ties going to the larger $\lambda$. The limit $\lambda\to\infty$, $A_\ell=0$ with $b_\ell$ the mean teacher state, which produces a constant forecast, takes the place of the hard cut as the explicit candidate. Penalizing $\|W A_\ell\|_F^2$ instead, i.e.\ plain ridge regression of the full model's head outputs on the states, gives similar results (Table~\ref{tab:timesfm3_ridge}, last aligned column), so the conclusions do not depend on the form of the penalty. The identity-anchored fit is reproduced exactly by the same code: its test MASE matches the submission to within $3.4\times10^{-5}$ relative at every depth and dataset, with the same selected $\kappa$ in all 280 fits.

\paragraph{Results.} The zero anchor changes the aligned TimesFM-3 pathway substantially (Table~\ref{tab:timesfm3_ridge}, Figure~\ref{fig:app_h3_depth_timesfm3}). From L1 to L14 its median test MASE increase stays between $+9.9\%$ and $+11.9\%$, against $+38.7\%$ to $+81.9\%$ with the identity anchor, and it is better on all 14 datasets at L1 to L8 and L10 to L15 (paired 95\% intervals). L0 becomes its worst depth: the zero-anchored pathway at L0 is worse than at every depth from L1 to L18 on 12 of the 14 datasets. Validation alone prefers the zero anchor in all 280 fits. Its median validation error is 1.9 to 3.2 times its training error, comparable to the iteratively fitted Chronos-2 and TiRex adapters (median 1.8 and 2.3 at the selected checkpoint).

The corrected aligned pathway also has a lower test MASE than the supervised readout on 10 to 12 of the 14 datasets at every depth from L1 (paired 95\% intervals). Both are affine maps of the same state, but the readout is fitted to the observed future with the pinball loss, the adapter to the full model's forecasts; we do not attribute the difference to a single cause. The exceptions are BOOM and Wind Farms, where validation selects $\kappa\ge10$, i.e.\ a nearly constant forecast, at L1 to L18; there, none of our affine fits reproduces the full model's forecasts on validation windows, and the aligned pathway remains far behind the readout ($+181\%$ and $+38\%$ at L5, against $+47\%$ and $+10\%$). On LOOP Seattle the readout is 8\% to 20\% better than the full model itself, which a fit to the full model's forecasts cannot follow.

\paragraph{Order of analysis and physical truncation.} We introduced the zero anchor after the test results of the identity-anchored fit had been computed, prompted by its anomalous L0 performance; both fits are reported here, and the main text uses the zero-anchored one. All TimesFM-3 accuracy numbers with the zero anchor are computed offline, from cached representations. The adapter has the same shape as before, so the measured latency is unchanged. We repeated the physical check of Appendix~\ref{app:latency} for the zero-anchored adapters at L0, L1, L5, L10, L15, L18, and L19 on all 14 datasets, 98 configurations evaluated on CPUs. Each adapter is refitted from the cached states with the same code, selects the same $\kappa$, and reproduces the offline numbers exactly; it is then spliced before the head of a copy of the public TimesFM-3 model with the blocks after L$\ell$ removed. In every configuration, the physically truncated model's mean test MASE agrees with the offline pathway to within $4.5\times10^{-5}$ relative, satisfying the rule of Appendix~\ref{app:latency}, and the unmodified model agrees with the cached full-model forecasts to within $4.2\times10^{-5}$. As for the original adapters, individual forecast values agree less tightly: the median difference is about $10^{-6}$ of the mean forecast magnitude, and the largest is 14\% on BOOM at L0 and at most 2.3\% elsewhere.

\begin{table}[t]
\centering
\caption{\textbf{TimesFM-3 aligned pathway with identity- and zero-anchored regularization.} Median test MASE increase over the full model across 14 datasets; $\times$ denotes ratios to the full-model MASE. Zero, $\|WA\|$: the zero-anchored fit with the penalty $\lambda\|W A_\ell\|_F^2$. $\le$5\%: datasets within 5\% of the full model (zero anchor). Zero $<$ identity and zero vs.\ readout: datasets where the paired 95\% interval of the ratio excludes one, better / worse. Speedup: end to end at batch size 1, from the physically truncated identity-anchored models; the zero-anchored adapter has the same shape.}
\vspace{4pt}
\label{tab:timesfm3_ridge}
\small
\setlength{\tabcolsep}{3.5pt}
\begin{tabular}{@{}lrrrrrcccr@{}}
\toprule
& & \multicolumn{3}{c}{Aligned pathway} & & & & & \\
\cmidrule(lr){3-5}
Cut & Hard cut & Identity & Zero & Zero, $\|WA\|$ & Readout & $\le$5\% & Zero $<$ id. & Zero vs.\ readout & Speedup \\
\midrule
L0 & $35\times$ & $+31.3\%$ & $+24.3\%$ & $+23.9\%$ & $+31.9\%$ & 0 & 8 & 8 / 2 & $1.38\times$ \\
L1 & $38\times$ & $+38.7\%$ & $+10.6\%$ & $+9.7\%$ & $+15.8\%$ & 2 & 14 & 10 / 2 & $1.37\times$ \\
L5 & $53\times$ & $+69.8\%$ & $+10.6\%$ & $+8.5\%$ & $+16.8\%$ & 4 & 14 & 10 / 4 & $1.24\times$ \\
L10 & $121\times$ & $+67.3\%$ & $+11.3\%$ & $+10.1\%$ & $+17.0\%$ & 2 & 14 & 10 / 3 & $1.19\times$ \\
L15 & $160\times$ & $+47.9\%$ & $+8.5\%$ & $+7.1\%$ & $+14.8\%$ & 3 & 14 & 10 / 3 & $1.08\times$ \\
L17 & $91\times$ & $+29.7\%$ & $+6.1\%$ & $+6.0\%$ & $+14.7\%$ & 5 & 13 & 11 / 3 & $1.07\times$ \\
L18 & $62\times$ & $+14.3\%$ & $+3.9\%$ & $+3.7\%$ & $+13.2\%$ & 9 & 13 & 11 / 3 & $1.03\times$ \\
L19 & $7\times$ & $+5.4\%$ & $+1.3\%$ & $+1.0\%$ & $+12.7\%$ & 11 & 13 & 12 / 1 & $0.99\times$ \\
\bottomrule
\end{tabular}
\end{table}

\begin{table}[t]
\centering
\caption{\textbf{What validation selected for the two TimesFM-3 adapters.} Medians over the 14 datasets at each depth, summarized over L1 to L17; L0 in parentheses. Effective parameters: the trace of the ridge smoother for each of the head's output directions, averaged with weights proportional to the full model's output variance in that direction (at most 1,280). Validation / training: distillation error of the selected fit on the validation windows relative to the training windows. Test $R^2$: against the full model's forecast in the head's output metric.}
\vspace{4pt}
\label{tab:timesfm3_ridge_fit}
\small
\setlength{\tabcolsep}{5pt}
\begin{tabular}{@{}lcc@{}}
\toprule
& Identity anchor (submission) & Zero anchor \\
\midrule
Selected $\kappa$ & $\le10^{-8}$ in 231 of 280 fits & $10^{-6}$ to $10^{-4}$ in 221 of 280 fits \\
Effective parameters & 1,017 to 1,280 (253) & 262 to 681 (76) \\
Validation / training error & 13 to 184 (1.9) & 1.9 to 3.2 (1.3) \\
Test $R^2$ & $-0.60$ to $0.53$ (0.12) & $0.64$ to $0.88$ (0.33) \\
Datasets with test $R^2<0$ & 4 to 11 (6) & 2 to 3 (3) \\
Lower validation error & 0 of 280 fits & 280 of 280 fits \\
\bottomrule
\end{tabular}
\end{table}
